\subsection{Every original DVR point and fixed-weight rigidity}
\label{sec:support-finite-points}

Retain the original coefficient DVR $O$, its finite residue field $k$
of odd characteristic $\ell$, the frame $F$ with all its original
\texttt{Ok} hypotheses, and the original normalized spaces $S_n$.
An \emph{original DVR lift} means a complete DVR $A$ with finite
residue field, a local $O$-algebra structure, and an original
normalized representation $\rho$ satisfying \texttt{IsLiftA}.
In particular, its residual representation is diagonal, with the
specified diagonal characters, and the representation and ordered
local characters are adically continuous.  These are the literal
point quantifiers, including residue extensions, in
\cite{support-finite-points-source}.

\begin{lemma}[A coefficient ring with the original residue field]
\label{lem:support-residue-pullback}
Every original DVR lift is the image of the original universal
representation under an $O$-algebra map $S_n.R\to A$.
\end{lemma}

\begin{proof}
Write $k'=k(A)$ and identify $k$ with its image in $k'$.  Set
\[
 C=A\mathbin{\times}_{k'}k
   =\{a\in A:\overline a\in k\}.
\]
This is a local $O$-algebra with maximal ideal $\mathfrak m_A$ and
residue field $k$.  Lifts of a $k$-basis of $k'$, together with $1$,
generate $A$ over $C$: subtract the corresponding combination with
coefficients lifted from $O$, leaving an element of
$\mathfrak m_A\subset C$.  For a uniformizer $\pi_A$,
\[
 \mathfrak m_C^r=\mathfrak m_A^r=\pi_A^rA\qquad(r\geq1),       \tag{FP1}
\]
since $\pi_A^rb$ is a product of $r$ elements of $C$.
Thus $C$ is closed and complete with its own adic topology.
For a nonzero proper ideal $J\subset C$, take $a\in J$ of least
$A$-valuation $r$.  The submodule
$a\mathfrak m_A=\pi_A^{r+1}A\subset J$ is finite over $C$, and
$J/(a\mathfrak m_A)$ is finite because $A/\pi_A^{r+1}A$ is finite.
Every ideal of $C$ is therefore finitely generated.  The field case,
if allowed in the DVR convention, simply gives $C=k$.

All entries of $\rho$, its inverse matrices, and the local character
values and inverses lie in $C$, since their prescribed residues lie
in $k$.  The normalized equations descend through the injection
$C\hookrightarrow A$, and (FP1) preserves their continuity.
The lift over $C$ is therefore in the original complete local
category with residue field $k$.  Its universal map
$S_n.R\to C$, followed by inclusion, is the required map.
\end{proof}

\begin{theorem}[The literal finite-point output]
\label{thm:support-finite-points}
For $n=2,3$, \texttt{AllGood} for the actual Hecke ideal implies
the original \texttt{GoodPt}: every original DVR lift satisfying
the original irreducibility test after any injective map to a field
has its full twisted determinant law and ordered local characters
factoring through the actual Hecke quotient over the same $A$.

For $n=3$, the original \texttt{OffFin} implies the literal
\texttt{FinPt}.  More precisely, let $A$ additionally be a
$k$-algebra with the original $O$--$k$--$A$ scalar tower, and assume
\[
 \psi_{v,i}(g)=1\quad
 (g\in I_v,\ \Psi(\operatorname{loc}_vg)=1).                 \tag{FP2}
\]
Then $\rho$ is globally upper triangular or globally lower triangular
in its original coordinates, its diagonal is
$(\Psi,1,\Psi^{-1})$, and for every $g$,
\[
 \det(X-\rho(g))=(X-\Psi(g))(X-1)(X-\Psi(g)^{-1}).           \tag{FP3}
\]
Consequently the original rank-two and rank-three \texttt{ModIn}
inputs give the complete original \texttt{AbsMain} output.
\end{theorem}

\begin{proof}
Use Lemma~\ref{lem:support-residue-pullback} and let
$P=\ker(A_n\to A)$.  It is prime.  For the first assertion,
irreducibility after the specified injective field map excludes
both triangular loci and, in rank three, the corner-plus-middle
splitting on $E$.  Thus $P$ is type I.  If $J$ is the original
Hecke ideal and $\Phi:B_{\rm det}\to B_n$ its compatible map,
\texttt{AllGood} gives
$\Phi(J)\subset P\cap B_n$.  The composite
$B_{\rm det}\to B_n\to A$ kills $J$ and carries the full determinant
law and ordered characters to the prescribed twisted ones by the
original compatibility.  This proves \texttt{GoodPt} integrally,
with every original $A$ and field-test quantifier.

For the second assertion, write $I\subset B_3$ for the actual
\texttt{OffFin} ideal.  Thus
\[
 \mathfrak m_{B_3}\subset\sqrt{I+\lambda_3},\qquad
 I\subset P'\cap B_3\quad\hbox{for every off-triangular prime }P'.
                                                               \tag{FP4}
\]
If $P$ were off triangular, the first ideal $I$ would vanish in $A$.
So would $\lambda_3$: its uniformizer vanishes through $O\to k\to A$,
and its original fixed-inertia generators vanish by (FP2).
Primality and (FP4) then kill $\mathfrak m_{B_3}$.
But off-triangularity supplies nonzero upper and lower entries
$u,v$ of degrees $a,-b$, where $a,b\in\{1,2\}$.
Their balanced product $u^bv^a$ belongs to $B_3$ and has zero
original residue, so lies in $\mathfrak m_{B_3}$.
Its image cannot vanish in the domain $A$.  This contradiction
proves global triangularity with the original coordinate ordering.

The diagonal entries $a_i$ are now continuous characters residually
$\Psi^{w_i}$.  The original ordered local equations give local
characteristic polynomial $\prod_i(X-\psi_{v,i})$; this follows from
the original \texttt{f3} element has residual $\Psi^2\ne1$, hence
three distinct residual eigenvalues and separating local projectors.
The resulting ordered flag over $\operatorname{Frac}A$ gives the identity,
which descends to $A$ by injection.  Under (FP2),
$(a_i(g)-1)^3=0$, hence $a_i(g)=1$, including when $\ell=3$.
The characters $b_i=a_i/\Psi^{w_i}$ have residue one.
For arbitrary inertia $g$, its power $g^{|k|-1}$ belongs to the
subgroup in (FP2), so $b_i(g)^{|k|-1}=1$ and $b_i(g)=1$.
Here the principal units of a characteristic-$\ell$ domain have
no nontrivial finite-order elements: $\ell$-power torsion is
excluded by Frobenius and prime-to-$\ell$ torsion by the simple
root at $1$.  The original polarization similarly makes every
$b_i$ trivial on conjugate inertia.

Modulo $\mathfrak m_A^N$, each $b_i$ has finite $\ell$-group image
and open kernel.  The literal original class-group exponent
condition \texttt{cl} gives one nonzero $h$ such that
$b_i(g)^h=1$ in every such quotient.  Krull intersection gives the
identity in $A$, and principal-unit torsion-freeness gives $b_i=1$.
The triangular determinant proves (FP3).

Apply the first assertion to the rank-two and rank-three
\texttt{AllGood} conclusions, and the second to their actual
rank-three \texttt{OffFin} ideal $I\cap I_3$.
These are precisely \texttt{GoodPt} in both ranks and
\texttt{FinPt} in rank three, the three outputs of \texttt{AbsMain}.
All original arithmetic \texttt{ModIn} hypotheses remain explicit.
\end{proof}

The triangularity assertion is about the original DVR domains.
It supplies neither a uniform nilpotence exponent over Artinian
rings nor an integral lattice-length bound.  The next arithmetic
interface is the original \texttt{GalLink} transfer to one fixed
finite-index Galois subgroup, followed by the two-sided finite-ring
criterion and the moving-family Selmer estimate.
